% System Design

\section{System Design}
\label{sec:design}

This section presents the design of \textsc{WasmSentinel}, focusing on its three key innovations: the WasmML detection engine, the hybrid detection architecture, and the zero-copy data channel.

\subsection{Architecture Overview}
\label{sec:design-overview}

\begin{figure*}[t]
\centering
\begin{tikzpicture}[scale=0.9, transform shape,
    box/.style={rectangle, draw, minimum height=0.8cm, align=center, font=\footnotesize},
    ebpf/.style={box, fill=blue!15, minimum width=1.6cm},
    host/.style={box, fill=green!15, minimum width=1.8cm},
    wasm/.style={box, fill=orange!15, minimum width=1.8cm},
    arrow/.style={->, >=stealth, thick}
]
    % Kernel space
    \node[ebpf] (syscall) at (0,0) {Syscall\\Tracing};
    \node[ebpf] (sched) at (1.9,0) {Process\\Lifecycle};
    \node[ebpf] (net) at (3.8,0) {Network\\Events};
    \node[ebpf] (file) at (5.7,0) {File\\Access};
    \node[ebpf] (ring) at (2.85,-1.2) {Ring Buffer};

    % User space
    \node[host] (reader) at (0,-3) {Event\\Reader};
    \node[host] (batcher) at (2.3,-3) {Batch\\Aggregator};
    \node[host] (engine) at (4.6,-3) {Wasm\\Engine};

    % Wasm modules
    \node[wasm] (stide) at (7.2,-2.5) {STIDE};
    \node[wasm] (mlp) at (7.2,-3.5) {MLP};
    \node[wasm] (hybrid) at (9.2,-3) {Hybrid\\Scheduler};

    % Arrows
    \draw[arrow] (syscall) -- (ring);
    \draw[arrow] (sched) -- (ring);
    \draw[arrow] (net) -- (ring);
    \draw[arrow] (file) -- (ring);
    \draw[arrow] (ring) -- (reader);
    \draw[arrow] (reader) -- (batcher);
    \draw[arrow] (batcher) -- (engine);
    \draw[arrow] (engine) -- (stide);
    \draw[arrow] (engine) -- (mlp);
    \draw[arrow] (stide) -- (hybrid);
    \draw[arrow] (mlp) -- (hybrid);

    % Labels
    \node[above=0.4cm of net, font=\small\bfseries] {Kernel Space};
    \node[above=0.3cm of batcher, font=\small\bfseries] {User Space};
    \node[above=0.3cm of hybrid, font=\small\bfseries] {Wasm Sandbox};

    % Boundary lines
    \draw[dashed, thick, gray] (-1.2,-0.6) -- (10,-0.6);
    \draw[dashed, thick, gray] (6,-4.2) -- (6,-1.8);
\end{tikzpicture}
\caption{\textsc{WasmSentinel} architecture. eBPF collects kernel events and emits them to a ring buffer. The host dispatches events to Wasm-sandboxed detectors.}
\label{fig:architecture}
\end{figure*}

Figure~\ref{fig:architecture} illustrates the overall architecture of \textsc{WasmSentinel}.
The system comprises three layers:

\textbf{Kernel Layer.}
eBPF programs attach to kernel tracepoints to capture security-relevant events including syscall invocations, process lifecycle events, network connections, and file accesses.
Events are emitted to a shared ring buffer for efficient transfer to user space.

\textbf{Host Layer.}
A Go-based host program reads events from the ring buffer, optionally batches them for efficiency, and dispatches them to the Wasm detection engine.
The host also manages Wasm module lifecycle, resource limits, and alert handling.

\textbf{Detection Layer.}
Detection logic executes within Wasm sandboxes, providing strong isolation guarantees.
Our hybrid architecture combines multiple detection strategies, each implemented as a Rust module compiled to Wasm.

\subsection{WasmML: ML in Wasm Sandbox}
\label{sec:design-wasmml}

The core innovation of \textsc{WasmSentinel} is \textit{WasmML}---the ability to run machine learning inference inside a Wasm sandbox while maintaining microsecond-level latency.

\subsubsection{Design Challenges}

Running ML inference in Wasm is non-trivial.
Standard frameworks like TensorFlow and PyTorch cannot run in Wasm---they depend on native libraries, GPU drivers, and dynamic loading, none of which are available.
Wasm also lacks SIMD instructions in most runtimes (though support is emerging), and linear memory must be pre-allocated rather than grown dynamically.
These constraints rule out simply porting existing ML code.

\subsubsection{Our Approach: Pure-Rust MLP}

Rather than porting heavyweight ML frameworks, we implement a lightweight MLP directly in Rust:

\begin{lstlisting}[style=rust, caption={Core MLP inference loop in Rust.}]
pub fn forward(&self, input: &[f32]) -> f32 {
    let mut x = input.to_vec();
    for (i, layer) in self.layers.iter().enumerate() {
        x = layer.weights.iter()
            .map(|row| {
                row.iter().zip(x.iter())
                   .map(|(w, x)| w * x)
                   .sum::<f32>()
            })
            .zip(layer.biases.iter())
            .map(|(sum, b)| {
                if i < self.layers.len() - 1 {
                    relu(sum + b)  // Hidden layers
                } else {
                    sigmoid(sum + b)  // Output
                }
            })
            .collect();
    }
    x[0]
}
\end{lstlisting}

This approach offers several advantages: zero external ML dependencies, predictable performance without garbage collection or JIT, a compact 180KB Wasm binary, and portable weights stored as JSON.

\subsubsection{Model Architecture}

Through systematic architecture search (Appendix~\ref{app:experiments}), we found that a 4-layer MLP achieves the best accuracy-latency trade-off:

\begin{table}[t]
\centering
\caption{MLP architecture for \textsc{WasmSentinel}.}
\label{tab:mlp-arch}
\begin{tabular}{lccc}
\toprule
\textbf{Layer} & \textbf{Input} & \textbf{Output} & \textbf{Activation} \\
\midrule
Input & 64 & 128 & ReLU \\
Hidden 1 & 128 & 64 & ReLU \\
Hidden 2 & 64 & 32 & ReLU \\
Hidden 3 & 32 & 16 & ReLU \\
Output & 16 & 1 & Sigmoid \\
\bottomrule
\end{tabular}
\end{table}

Raw feature extraction produces a 340-dimensional vector (Appendix~\ref{app:features}), which we compress to 64 dimensions via feature hashing before MLP input.
The compressed vector preserves syscall categories, dangerous syscall indicators, privilege context, and n-gram statistics while reducing inference cost.

Total model parameters: 19,201 (75KB as float32).

\subsection{Hybrid Detection Architecture}
\label{sec:design-hybrid}

\begin{figure}[t]
\centering
\begin{tikzpicture}[
    box/.style={rectangle, draw, minimum height=0.7cm, align=center, font=\footnotesize},
    decision/.style={diamond, draw, aspect=2.5, font=\tiny, inner sep=1pt},
    arrow/.style={->, >=stealth}
]
    \node[box, fill=blue!10] (event) at (0,0) {Event};
    \node[box, fill=green!10] (stide) at (2,0) {STIDE};
    \node[decision, fill=yellow!10] (check) at (4.2,0) {Anomaly?};
    \node[box, fill=orange!10] (mlp) at (6.4,0.8) {MLP};
    \node[box, fill=green!10] (normal) at (6.4,-0.8) {Normal};
    \node[decision, fill=yellow!10] (mlpcheck) at (8.6,0.8) {Threat?};
    \node[box, fill=red!10] (alert) at (10.5,0.8) {Alert};
    \node[box, fill=green!10] (benign) at (10.5,0) {Benign};

    \draw[arrow] (event) -- (stide);
    \draw[arrow] (stide) -- (check);
    \draw[arrow] (check) -- node[above, font=\tiny] {Yes} (mlp);
    \draw[arrow] (check) -- node[right, font=\tiny] {No} (normal);
    \draw[arrow] (mlp) -- (mlpcheck);
    \draw[arrow] (mlpcheck) -- node[above, font=\tiny] {Yes} (alert);
    \draw[arrow] (mlpcheck) -- node[right, font=\tiny] {No} (benign);
\end{tikzpicture}
\caption{Hybrid detection flow: STIDE filters normal events, suspicious events escalate to MLP.}
\label{fig:hybrid-flow}
\end{figure}

Pure ML-based detection incurs inference cost for every event.
To reduce overhead, we employ a \textit{hybrid detection architecture} that combines STIDE-based fast-path filtering with MLP-based deep analysis.

\subsubsection{STIDE Fast Path}

STIDE (Sequence Time-Delay Embedding)~\cite{stide} detects anomalies through n-gram analysis of syscall sequences.
During training, STIDE builds a database of ``normal'' n-grams observed in benign behavior.
At runtime, incoming syscall sequences are checked against this database; sequences containing novel n-grams are flagged as anomalous.

Our implementation uses 6-grams (empirically optimal for syscalls), FNV-1a hashing for speed and collision resistance, and stores up to 100,000 unique n-grams.

STIDE overhead is sub-microsecond per event (median 180ns over 100K events), dominated by hash computation.

\subsubsection{Escalation Policy}

The key insight behind hybrid detection is that most events are benign and can be classified quickly.
STIDE first computes an anomaly score based on the ratio of novel n-grams in the recent syscall window.
If this score falls below threshold $\theta_1$ (default: 0.3), the event is classified as benign without invoking ML---this handles 80--90\% of events in typical workloads.

Events exceeding $\theta_1$ are escalated to the MLP, which extracts features and returns a threat probability.
If this probability exceeds $\theta_2$ (default: 0.7), an alert is generated.
The two thresholds provide tunable trade-offs: lowering $\theta_1$ increases recall but also MLP invocations; raising $\theta_2$ reduces false positives but may miss subtle attacks.

\subsection{Zero-Copy Data Channel}
\label{sec:design-zerocopy}

Efficient data transfer from kernel to Wasm is critical for overall performance.
Traditional approaches incur multiple memory copies and serialization overhead.

\subsubsection{Traditional Approach}

Typical pipelines serialize kernel events into JSON, copy to userspace heap, then re-copy into Wasm memory---incurring 4 copies plus parsing (5--50$\mu$s per event).

\subsubsection{Our Approach: Direct Memory Mapping}

\textsc{WasmSentinel} reduces this to a single copy by mapping the ring buffer directly into Wasm linear memory.
The eBPF program writes events in a fixed binary format matching the Wasm-side struct layout.
The host maps a read-only view of the ring buffer region into the Wasm module's linear memory, allowing the detection code to read event fields at known offsets without any parsing or copying.

\begin{lstlisting}[style=rust, caption={Zero-copy event access in Wasm.}]
// Event structure matches kernel layout
#[repr(C)]
struct RawEvent {
    timestamp: u64,
    pid: u32,
    ppid: u32,
    uid: u32,
    event_type: u32,
    syscall_nr: u32,
    // ... additional fields
}

fn read_event(ptr: *const u8) -> &RawEvent {
    unsafe { &*(ptr as *const RawEvent) }
}
\end{lstlisting}

This approach reduces per-event overhead to under 500ns.

\subsubsection{Safety Considerations}

Direct memory mapping raises safety concerns that we address through read-only access (Wasm cannot modify the ring buffer), bounds checking against buffer boundaries, format validation via magic numbers and checksums, and double-fetch protection by copying events to stack before multi-field access.
